\newif\ifuseseminar
\useseminartrue
\input{../../latex_common/header.tex.frag}

\title{Prova de Física Matemática II -- EDO (Recuperação)}

\begin{document}

\begin{enumerate}

	\item Considere a equação diferencial
	      $$
		      x^2 y'' + (2x+x^2)y' - 2y = 0.
	      $$
	      \begin{enumerate}
		      \item Classifique o ponto $x_0 = 0$.
		      \item Reescreva a equação na forma apropriada para o método de Frobenius,
		            determine as funções analíticas $p(x)$ e $q(x)$ em torno de
		            $x_0 = 0$ e obtenha o polinômio indicial.
		      \item Determine a relação de recorrência para os coeficientes da série.
	      \end{enumerate}

	\item Considere a equação diferencial
	      $$
		      x^2 y'' + xy' + (x^2 - 4)y = 0.
	      $$
	      \begin{enumerate}
		      \item Mostre que a equação possui um ponto singular regular em
		            $x_0 = 0$.
		      \item Encontre a solução correspondente ao maior índice da equação
		            indicial usando o método de Frobenius (polinômio indicial e
		            relação de recorrência).
		      \item Classifique o tipo da solução linearmente independente (Caso 1,
		            Caso 2-A, Caso 2-B-i ou Caso 2-B-ii) e discuta o aparecimento
		            de termos logarítmicos.
	      \end{enumerate}

	\item Considere a equação diferencial
	      $$
		      xy'' + (2-x)y' + \lambda y = 0.
	      $$
	      \begin{enumerate}
		      \item Resolva a equação em torno de $x_0 = 0$ por séries de potências
		            e obtenha a relação de recorrência.
		      \item Determine para quais valores de $\lambda$ a solução se reduz a
		            um polinômio.
	      \end{enumerate}

	\item Considere a equação
	      $$
		      (1-x^2)y'' - xy' + \nu y = 0.
	      $$
	      \begin{enumerate}
		      \item Mostre que $x=\pm1$ são pontos singulares regulares.
		      \item Reescreva a equação na forma apropriada para o método de
		            Frobenius, determine as funções analíticas $p(x)$ e $q(x)$ em
		            torno de $x_0=-1$ e encontre os expoentes indiciais.
	      \end{enumerate}

\end{enumerate}

\end{document}